
\medskip
\phantomsection\label{prop:sga3-VIB-7.6}%
\noindent\textbf{Proposition 7.6.---}
\emph{Let \(G\) be a \(k\)-group locally of finite type, \(X\) a
geometrically reduced \(k\)-scheme locally of finite type, and
\(f:X\to G\) a \(k\)-morphism. For the functor \(\Gamma_G(f)\) to
be the sheaf associated, for the \((\mathrm{fppf})\) topology
\textup{(}or the \((\mathrm{fpqc})\) topology\textup{)}, with the
group subpresheaf of \(G\) generated by the image of the morphism
of functors \(f\), it is necessary and sufficient that}
\phantomsection\label{form:sga3-VIB-7.6-union}%
\[
  \bigcup_{n\geq1}f^n(X^n)=\Gamma_G(f).
\]
\emph{When \(G\) is of finite type over \(k\), this condition
already implies that there exists an integer \(n\) such that
\(f^n:X^n\to\Gamma_G(f)\) is a covering for the
\((\mathrm{fpqc})\) topology \textup{(}and a fortiori
surjective\textup{)}.}\footnote{Translator's note: this last assertion is not correct as stated, and the end of the proof, which
derives it from EGA~IV\(_1\), 1.9.9, only shows that finitely many of the \(f^n(X^n)\) cover
\(\Gamma_G(f)\). For example, if \(G\) is the constant \(k\)-group \((\mathbb Z/2\mathbb Z)_k\),
\(X=\operatorname{Spec}k\) and \(f\) is the point \(1\), then \(\Gamma_G(f)=G\) and the condition holds, but
each \(f^n(X^n)\) is the single point \(n\bmod 2\). The argument does show that, for some \(N\), the morphism
\(\coprod_{n\leq N}X^n\to\Gamma_G(f)\) defined by the \(f^n\) is a covering (apply 7.5 to it, as in the
proof); and when \(f(X)\) contains the unit element, the \(f^n(X^n)\) increase with \(n\), some \(f^n\) is
surjective, and then \(f^{2n}\) is a covering by 7.5. The Gille--Polo re-edition (2011),
Exposé~VI\textsubscript{B}, keeps the assertion, in the general form of the Addition below.}

Put
\[
  X^\infty=\coprod_{n\geq1}X^n,
\]
and let \(f^\infty:X^\infty\to\Gamma_G(f)\) be the morphism
whose restriction to each \(X^n\) is the morphism
\(X^n\to\Gamma_G(f)\) deduced from \(f^n\). The sheaf considered
in the statement is then the image sheaf of \(X^\infty\) under
\(f^\infty\); hence (Exposé~IV,
\SGARefLink{sga3:iv:proposition:4:4:3}{4.4.3}), saying that the
functor \(\Gamma_G(f)\) is the sheaf considered amounts to saying
that \(f^\infty\) is a covering. Let us show that \(f^\infty\) is
a covering if and only if it is surjective. The condition is
obviously necessary; conversely, if \(f^\infty\) is surjective,
then
\phantomsection\label{form:sga3-VIB-7.6-finfty-product}%
\[
  \mu\circ
  \bigl(f^\infty\times_k f^\infty\bigr):
  X^\infty\times_k X^\infty\longrightarrow\Gamma_G(f)
\]
is a covering (\hyperref[lem:sga3-VIB-7.5]{7.5}); now, since
\(X^n\times_kX^m\) is canonically isomorphic to \(X^{n+m}\), and
since under this isomorphism \(\mu\circ(f^n\times_kf^m)\)
corresponds to \(f^{n+m}\), it is clear that
\(\mu\circ(f^\infty\times_kf^\infty)\) factors through
\(f^\infty\), so that \(f^\infty\) is a covering. Finally, saying
that \(f^\infty\) is surjective amounts to saying that
\(\bigcup_{n\geq1}f^n(X^n)=\Gamma_G(f)\). If \(G\), and hence
\(\Gamma_G(f)\), is of finite type, one concludes that there
exists \(n\) with \(f^n(X^n)=\Gamma_G(f)\), because the
\(f^n(X^n)\) are ind-constructible subsets of \(\Gamma_G(f)\)
(EGA~IV\(_1\), 1.8.4) with union \(\Gamma_G(f)\), and one applies
EGA~IV\(_1\), 1.9.9. This completes the proof of 7.6.

\medskip\noindent
\emph{Addition} (after the Gille--Polo re-edition (2011), Exposé~VI\textsubscript{B}; our
formulation). The re-edition states 7.6 in a relative form, noting that this form is used implicitly in
the proof of \hyperref[prop:sga3-VIB-10.12]{10.12}. Let \(S\) be a scheme, \(H\) an \(S\)-group locally of
finite presentation over \(S\) with reduced fibers, \(X\) an \(S\)-scheme flat and locally of finite
presentation over \(S\), and \(f\colon X\to H\) an \(S\)-morphism. Define \(X^n\) and \(f^n\colon X^n\to H\)
as in 7.1~(ii), with products over \(S\), put \(X^\infty=\coprod_{n\geq1}X^n\), and let
\(f^\infty\colon X^\infty\to H\) be the morphism induced by the \(f^n\). For every \(S\)-scheme \(S'\), the
subgroup of \(H(S')\) generated by \(f(X(S'))\) consists of the products
\(f(x_1)^{\pm1}\cdots f(x_n)^{\pm1}\), \(x_i\in X(S')\); so this group subpresheaf of \(H\) and the image
presheaf of \(f^\infty\) coincide on connected \(S'\), and have the same associated sheaf. \emph{The following
conditions are equivalent: \textup{(i)} \(H\) represents the \((\mathrm{fppf})\) sheaf associated with this
presheaf; \textup{(ii)} \(f^\infty\) is a covering for the \((\mathrm{fppf})\) topology; \textup{(iii)}
\(f^\infty\) is surjective, that is, \(H=\bigcup_{n\geq1}f^n(X^n)\). If, moreover, \(H\) is quasi-compact,
then for some \(N\) the morphism \(\coprod_{n\leq N}X^n\to H\) defined by the \(f^n\) is an
\((\mathrm{fppf})\) covering.} (The re-edition asserts this for a single \(f^n\), which fails in general:
see the note to the statement of 7.6.) Indeed, (i) amounts to (ii) (Exposé~IV,
\SGARefLink{sga3:iv:proposition:4:4:3}{4.4.3}), and (ii) implies (iii). Conversely, assume (iii); let
\(s\in S\), let \(\eta\) be a maximal point of \(H_s\), and let \(x\in X^\infty\) with \(f^\infty(x)=\eta\). As
\(H_s\) is reduced, its local ring at \(\eta\) is a field, so \(f^\infty_s\) is flat at \(x\); since
\(X^\infty\) is flat and locally of finite presentation over \(S\) and \(H\) is locally of finite
presentation over \(S\), \(f^\infty\) is flat at \(x\) (EGA~IV\(_3\), 11.3.10). Lemma~7.5 therefore applies to
\(f^\infty\): the morphism \(\mu\circ(f^\infty\times_Sf^\infty)\) is an \((\mathrm{fppf})\) covering, and it
factors through \(f^\infty\), as in the proof of 7.6, so \(f^\infty\) is a covering. If \(H\) is
quasi-compact, each \(f^n\) is locally of finite presentation (EGA~IV\(_1\), 1.4.3~(v)), so the
\(f^n(X^n)\) are ind-constructible (EGA~IV\(_1\), 1.9.5~(viii)), and by EGA~IV\(_1\), 1.9.9, finitely many of
them, say those with \(n\leq M\), cover \(H\); the same argument, applied to
\(\coprod_{n\leq M}X^n\to H\), shows that \(\coprod_{n\leq 2M}X^n\to H\) is a covering. For \(S=\operatorname{Spec}k\)
and \(H=\Gamma_G(f)\), with \(G\), \(X\) and \(f\) as in 7.6, one recovers the first assertion of 7.6, and the
second in the corrected form of the note to its statement.
